\section{Statistical physics}\label{sec:stat_phys}
\subsection{Brownian motion}\label{brown_mot}
Stochastic calculus has been fundamental in understanding complex systems: separating the dynamics of macroscopic observables into deterministic terms and stochastic terms reduces high-dimensional systems to a small number of variables of interest. Across physics, economics, geology, climate science, machine learning, and neuroscience, stochastic processes are often the most natural, tractable, and accurate modelling choice when a system is too complex to model in full. The important concept here is that of the introduction of some central limit theorem, where the sum of very many small interactions which are understood microscopically can produce certain well-defined stochastic patterns on macroscopic scales of interest. Because sums of classical random variables tend toward Gaussian, the most ubiquitous stochastic process is a Wiener process $W_t$, equivalently called Brownian motion, which has stationary and independent Gaussian increments across all non-overlapping time intervals. The Wiener process is a very simple and analytically elegant model of the coarse-grained effect of very many small finite-variance interacting influences. A Wiener process is defined only by its timescale, with a variance over trajectories after time $t$ of $t$ itself. Wiener processes yield time series which are almost surely continuous but nowhere differentiable. The Wiener process is understood as a model of the long-time dynamics of a variable effected by very many small uncorrelated finite variance influences \cite{hassler_stochastic_2016}. Classical examples of processes modelled by the Wiener process are: the movement of a macroscopic particle in a non-turbulent fluid \cite{hassler_stochastic_2016}, the gravitational influence of an orbiting thermal collection of stars on supermassive black holes \cite{Chandrasekhar1943, Merritt2007a}, and (arguably) the movement of low-volatility stock prices \cite{cunchala2024basicoverviewvariousstochastic}. 

However, there are well-studied mechanisms through which complex real-world stochastic processes diverge from the Wiener process. Only recently have the universality of stochastic processes asides from the ordinary Wiener process been understood. Of interest to us is the case where the microscopic effects become correlated and yield heavy-tailed statistics. Correlated steps, long-time memory effects, and microscopic dynamics can result in a stochastic process which does not limit to the Wiener process \cite{Brevitt2025, Klages2013}. Of interest to us is the fractional and superdiffusive Lévy flight, which is a result of spatial non-locality in an interacting network. For instance, previous studies have demonstrated a neural circuit network at criticality resulting in diverging correlation length of activity, and the centre of mass of neural firing is subject to the large jumps. Such large jumps in a stochastic time series are modelled according to Lévy stable motion, where the increments are drawn independently from a heavy-tailed distribution. This section describes the analytics and behaviour of such Lévy processes for our current study.


\subsection{Anomalous diffusion}\label{subsec:anom_diff}
A key motivation for the development for stochastic modelling of phenomena in nature was the study of diffusion. In 1885, the German physiologist Adolf Fick was the first to describe the universality of normal diffusion in the random motion of a large number of particles, producing a gradient of concentration in a fluid \cite{klafter_anomalous_2005}. Fick records the diffusion equation, which is identical to the heat equation applied to a probability distribution, as
\begin{equation}
    \frac{\partial p(x,t)}{\partial t} = D \frac{\partial ^2 p(x,t)}{\partial x ^2}
\end{equation}
where $D$ is a scaling constant. Fick shows that this gives a mean squared displacement (MSD) $\langle x(t)^2 \rangle = 2Dt$. This MSD linear in time defines normal diffusion. Such normal diffusion is found in many simple random processes, such as the paradigm case of the Brownian motion of dust in a fluid studied by Albert Einstein. Both Albert Einstein and the British statistician Karl Pearson independently developed the continuous time random walk (CTRW) with finite-variance step size and step time distributions, robustly producing this normal diffusion. Increasingly, physical systems have been which diffuse nonlinearly in time, a situation known as anomalous diffusion, with an MSD scaling as:
\begin{equation}
    \langle x(t) ^2 \rangle \propto t^{\beta}, \ \ \beta \neq 1
\end{equation}
Such anomalous diffusion yields two cases, subdiffusion $0 < \beta < 1$ where diffusion scales slower than normal, and superdiffusion $1 < \beta < 2$ where diffusion scales faster than normal. The case of $\beta = 2$ is known as ballistic diffusion, as it is found in the deterministic case of constant velocity $x = vt$. Anomalous subdiffusion has been shown to result from heavy-tailed wait-times between interactions \cite{klafter_anomalous_2005}, such as in trapped fluid flow \cite{Kirchner2000-kf}, controlled diffusion through cell membranes \cite{klafter_anomalous_2005}, and protein mixing \cite{Kou2004-kz}. \\

\subsubsection*{Anomalous superdiffusion}
More importantly to this thesis, superdiffusion has been studied in the contexts of complex exploratory biological systems. Superdiffusion has a strong link to the problem of a blind search for scarce resources, however the true computational benefit of superdiffusion has become a subject of debate in the literature. On one side, the Lévy foraging hypothesis posits that the standard superdiffusive CTRW, the Lévy walk, is an optimal strategy when searching for sparse and randomly distributed resources. This result was first reported in a letter to the journal Nature in 1999 by \textcite{viswanathan_optimizing_1999}. However, recent physicists and biologists have questioned if this is instead an artefact of improper modelling, an argument led by described most strongly this year by \textcite{pyke_levy_2026}. In the original paper proposing the Lévy foraging hypothesis a Lévy walk with superdiffusion $\langle x(t) ^2 \rangle \propto \frac{t^2}{\ln(t)}$ on the transition to ballistic diffusion was found to be optimal in a search for sparse resources \cite{viswanathan_optimizing_1999}. Further, the optimality of the superdiffusive Lévy walk has been reproduced in diverse functional search models \cite{humphries_optimal_2014, campeau_evolutionary_2022, sims_optimal_2019}, in a population of diverse searching agents with a range of superdiffusive exponents \cite{clementi_search_2021}, and have been reported in the empirical search paths of biological organisms from the flights of drosophila flies \cite{sims_optimal_2019, reynolds_free-flight_2007}, albatross' searches \cite{focardi_levy_2014, humphries_new_2013}, and strongly across many marine predators \cite{sims_scaling_2008, humphries_environmental_2010}. Despite these analyses, advantages of Lévy walks have been shown to be model-dependent, with the optimal performance of Lévy walks requiring biologically inaccurate depictions of search paths \cite{pyke_levy_2026, gautestad_levy_2013, pyke_understanding_2015}. These arguments marks a split in the treatment of Lévy walks, with groups studying them as informative baseline functional models of abstract search \cite{viswanathan_optimizing_1999, gautestad_levy_2013, humphries_optimal_2014, clementi_search_2021}, and others led by Graham Pyke presenting this as a misrepresentation of true intelligent organisms \cite{pyke_levy_2026, pyke_understanding_2015, benhamou_ultimate_2015}. This thesis provides a biologically realistic and well-motivated superdiffusive search for research into a case for optimal Lévy motion. See \textcite{reynolds_liberating_2015} for an in-depth discussion of a richer view of biological Lévy walk superdiffusion.
% Simulating trajectories from stochastic processes such as Wiener processes and Lévy flights are computationally direct, it is often useful to understand the dynamics of the distribution of trajectories after some time $t$. The dynamics of the distribution of trajectories driven by a stochastic differential equation is defined through the Fokker-Planck Equation (FPE), also known as the Kolmogorov Forward Equation. The paradigmatic FPE is given for the Langevin Equation:
% \begin{equation}
% 	dX_t = \mu(X_t,t)dt + \sigma(X_t,t) dW_t
% \end{equation}
% giving the partial differential equation for the distribution $p(x,t)$ of $X_t$:
% \begin{equation}
% 	\frac{\partial p(x,t)}{\partial t} = - \frac{\partial}{\partial x} [ \mu(x, t)p(x,t) ] + \frac{1}{2} \frac{\partial^2}{\partial x^2}[\sigma^2(x,t)p(x,t)]
% \end{equation}
% In the case where $\mu(x,t) = 0$ and $\sigma(x,t) = \sigma$, the Langevin FPE simplifies into the standard heat equation, with diffusion of the probability mass with a rate defined by the volatility parameter $\sigma$:
% \begin{equation}
% 	\frac{\partial p(x,t)}{\partial t} = \frac{\sigma^2}{2} \nabla^2 p(x,t)
% \end{equation}
% The heat equation is very well studied, and given an initial condition concentrated at the point $(x,t) = (x_0,0)$, the distribution at any time $t>0$ is exactly a Gaussian defined by:

% \begin{equation}
% 	p(x,t) = \frac{1}{(2 \pi \sigma^2 t)^{1/2}} \exp \left( - \frac{|x|^2}{2 \sigma^2 t} \right)
% \end{equation}

% We see from this that the variance of such a distribution will scale as $\textbf{Var}_x[p(x,t)] \propto t$

\subsection{Lévy stable motion}\label{subsec:levy_motion}
To mathematically capture the intermittent, scale-free nature of visual search, we turn to the \textit{Lévy flight}: a superdiffusive stochastic process that allows very strong analysis, due to it's very foundational mathematical definition. Here, the Lévy flight $L_T^{\alpha}$ is a stochastic process simply defined through it's increments being drawn from a heavy tailed symmetric stable distribution (a $\mathcal{S}\alpha\mathcal{S}$ distribution). $\mathcal{S}\alpha\mathcal{S}$ distributions describe the limit of summation of i.i.d infinite-variance variables, as a result of the generalised central limit theorem, making them ideal to describe a universality class of superdiffusive processes. The $\alpha$ value defining a particular $\mathcal{S}\alpha\mathcal{S}$ distribution is a result of the asymptotic power-law tail of the density function, with the probability density function following $p(x) \propto |x|^{-1-\alpha}$. Thus, the Lévy flight $L_t^{\alpha}$ after time $t$ is defined as follows, to outline the fractional scaling properties and alpha-stable statistics
\begin{equation}
	L_t \stackrel{d}{=} t^{1/\alpha} L_1, \ \ L_1 \sim \mathcal{S}\alpha\mathcal{S}
\end{equation}
where the parameter $\alpha$ defines the heaviness of the tails of resulting statistics, and where $\stackrel{d}{=}$ denotes equal in distribution. In the limit $\alpha \rightarrow 2$, the Lévy flight reduces to a scaled Brownian Motion $L_t^{\alpha \rightarrow 2} = \frac{1}{\sqrt{2}}W_t$. If $1 < \alpha < 2$ the variance of $L_t^{\alpha}$ diverges and has mean 0, whereas if $0 < \alpha \leq 1$, then both the variance and mean diverge. In analysing the distribution of a Lévy flight, one must employ fractional calculus in the form of the fractional Fokker-Planck equation. The details of such an analysis are beyond this thesis, however the fractional Fokker-Planck equation is the tool used to understand the superdiffusion,  nonlocality, and scale-free behaviour of trajectories of Lévy flights (see the appendices to \cite{qi_fractional_2022}). Many of the simple analytic properties of the Wiener process are no longer true of a Lévy flight. Importantly, trajectories of a Lévy flight are almost surely no longer continuous. Instead, such trajectories are called \textit{càdlàg} (continue à droite, limite á gauche) with respect to time, which describes functions which are right continuous with well-defined left limits \cite{van_den_Bosch}. Thus, we say that since all such trajectories produced from a Lévy flight are càdlàg, the Lévy flight itself is called a càdlàg process. Càdlàg processes have been widely studied in the context of mathematics, economics, and physics, and their properties are well understood \cite{soulier2022-ge, Borkar1995}. It has been shown that Lévy flights are semimartingales, and are well interpreted in integration \cite{van_den_Bosch}. Importantly, it follows from the theorem of the Lévy-Itô decomposition that a Lévy flight can be decomposed as the sum of two processes \cite{Kyprianou2014}:
\begin{enumerate}
	\item A compound Poisson process of large jumps with $|\Delta x| > 1$. Such jumps are exponentially distributed in time, and there are almost surely a finite number of such jumps in any finite time period.
	\item A square-integrable martingale (mean 0) jump process with $|\Delta x| < 1$. This process is thus finite-variance, and is thus subject to the central limit theorem, meaning that under temporal coarse graining it approaches a Wiener process.
\end{enumerate}
This division is useful in analysing the process in the context of nonlocal stochastic exploration in Lévy flights, which is modelled by the compound Poisson process of large jumps $|\Delta x| > 1$, and the local exploitation dynamics modelled by the square-integrable martingale of small jumps $|\Delta x| < 1$.